\begin{thebibliography}{17}
\bibitem{r1} Y. Zhong and X. Chen, ``An FFT Twofold Subspace-Based Optimization Method for Solving Electromagnetic Inverse Scattering Problems,'' \emph{IEEE Trans. Antennas Propag.}, vol. 59, no. 3, pp. 914--927, 2011, doi: 10.1109/TAP.2010.2103027.
\bibitem{r2} X. Ye and X. Chen, ``Subspace-Based Distorted-Born Iterative Method for Solving Inverse Scattering Problems,'' \emph{IEEE Trans. Antennas Propag.}, vol. 65, no. 12, pp. 7224--7232, 2017, doi: 10.1109/TAP.2017.2766658.
\bibitem{r3} X. Chen, \emph{Computational Methods for Electromagnetic Inverse Scattering}. Wiley--IEEE Press, 2018, Secs. 6.3.1 and 6.4.2.
\bibitem{r4} E. de Sturler, S. Gugercin, M. Kilmer, S. Chaturantabut, C. Beattie, and M. O'Connell, ``Nonlinear Parametric Inversion Using Interpolatory Model Reduction,'' \emph{SIAM J. Sci. Comput.}, vol. 37, no. 3, pp. B495--B517, 2015, doi: 10.1137/130946320. Author preprint: arXiv:1311.0922.
\bibitem{r5} M. Kartmann, T. Keil, M. Ohlberger, S. Volkwein, and B. Kaltenbacher, ``Adaptive Reduced Basis Trust Region Methods for Parameter Identification Problems,'' arXiv:2309.07627v3, 2024.
\bibitem{r6} T. Bui-Thanh, K. Willcox, O. Ghattas, and B. van Bloemen Waanders, ``Goal-Oriented, Model-Constrained Optimization for Reduction of Large-Scale Systems,'' \emph{J. Comput. Phys.}, vol. 224, no. 2, pp. 880--896, 2007, doi: 10.1016/j.jcp.2006.10.026.
\bibitem{r7} G. Meurant and P. Tichy, ``Error Norm Estimates for the Block Conjugate Gradient Algorithm,'' arXiv:2502.14979, 2025.
\bibitem{r8} A. Gaul, M. H. Gutknecht, J. Liesen, and R. Nabben, ``A Framework for Deflated and Augmented Krylov Subspace Methods,'' \emph{SIAM J. Matrix Anal. Appl.}, vol. 34, no. 2, pp. 495--518, 2013, doi: 10.1137/110820713.
\bibitem{r9} K. Smetana and A. T. Patera, ``Optimal Local Approximation Spaces for Component-Based Static Condensation Procedures,'' \emph{SIAM J. Sci. Comput.}, vol. 38, no. 5, pp. A3318--A3356, 2016, doi: 10.1137/15M1009603.
\bibitem{r10} K. Smetana, O. Zahm, and A. T. Patera, ``Randomized Residual-Based Error Estimators for Parametrized Equations,'' arXiv:1807.10489v2, 2018.
\bibitem{r11} B. T. Draine and P. J. Flatau, ``Discrete-Dipole Approximation for Scattering Calculations,'' \emph{J. Opt. Soc. Am. A}, vol. 11, no. 4, pp. 1491--1499, 1994. Publisher record.
\bibitem{r12} T. Cui et al., ``Likelihood-informed dimension reduction for nonlinear inverse problems,'' arXiv:1403.4680, 2014.
\bibitem{r13} A. Spantini et al., ``Goal-oriented optimal approximations of Bayesian linear inverse problems,'' arXiv:1607.01881, 2016.
\bibitem{r14} FAEMIS Project Team, ``FAEMIS Experiment Roadmap v5: Fresnel-Aware Electromagnetic Inverse Scattering,'' unpublished internal technical note, 2026. Used here only to define the project's modeling scope.
\bibitem{r15} C. F. Bohren and D. R. Huffman, \emph{Absorption and Scattering of Light by Small Particles}. Wiley-VCH, 1998, doi: 10.1002/9783527618156.
\bibitem{r16} X. Chen, Y. Zhong, and K. Agarwal, ``Subspace Methods for Solving Electromagnetic Inverse Scattering Problems,'' \emph{Methods and Applications of Analysis}, vol. 17, no. 4, pp. 407--432, 2010.

\bibitem{r17} Y. Zhong and X. Chen, ``Twofold Subspace-Based Optimization Method for Solving Inverse Scattering Problems,'' \emph{Inverse Problems}, vol. 25, art. 085003, 2009.
\end{thebibliography}